% ============================================
% DOC-02
% Gerado de molambudos.md — build_miolo.py
% ============================================
\fragtitle{doc}{D}{DOC}{02}{Prontuário do Paciente 1.260}
\markboth{DOC-02 --- Prontuário do Paciente 1.260}{}

\noindent \textit{Hospital Colônia de Barbacena --- Ficha de Admissão}

\noindent \textbf{Hospital Colônia de Barbacena --- Ficha de Admissão}
\textit{Reconstituição ficcional de ficha de admissão; não corresponde a documento arquivístico autêntico.}
\textit{Composição narrativa baseada em contexto histórico, sem chancela do Arquivo Público Mineiro.}

\vspace{0.5em}\noindent\rule{\textwidth}{0.5pt}\vspace{0.5em}

\noindent \textbf{Nº de Registro:} 1.260
\textbf{Nome:} Joaquim Antônio Correia
\textbf{Data de Nascimento:} ~1907 (estimado)
\textbf{Naturalidade:} Quixeramobim, Ceará
\textbf{Filiação:} Pai: Antônio Correia (falecido --- inanição, 1915); Mãe: Maria do Rosário Correia (desaparecida --- realocação forçada, 1915)
\textbf{Procedência:} Retirante saído de Senador Pompeu em 1915 e confinado no Campo do Alagadiço, próximo a Fortaleza. Transferido em 1917 via Santa Casa de Misericórdia de Belo Horizonte.

\vspace{0.5em}\noindent\rule{\textwidth}{0.5pt}\vspace{0.5em}

\noindent \textbf{Exame de Admissão --- Dr. Nelson Fontes, 12 de novembro de 1917:}

\noindent \textit{\textquotedblleft{}Paciente apresenta-se em estado de caquexia severa. Peso estimado: 20 kg (esperado para idade: 28-30 kg). Desnutrição proteico-calórica evidente. Edema nos membros inferiores. Abdome distendido. Dentição comprometida (cáries múltiplas, perda de 3 dentes).}

\noindent \textit{Sinal mais notável: heterocromia da íris direita, que se apresenta em tom amarelo-ouro, ausente na esquerda. Paciente alega que 'um bicho' mudou a cor. Provável sequela infecciosa ou traumática não identificada.}

\noindent \textit{Estado mental: pacato, cooperativo, mas com olhar fixo e vago. Responde perguntas com monossílabos. Quando questionado sobre a família, declarou: 'Morreram. Todos.' Sem sinais de alteração de sensório. Orientado no tempo e espaço.}

\noindent \textit{Avaliação preliminar: Esquizofrenia hebefrênica (demência precoce). Recomenda-se isolamento e observação.\textquotedblright{}}

\vspace{0.5em}\noindent\rule{\textwidth}{0.5pt}\vspace{0.5em}

\noindent \textbf{Histórico de Incidentes:}

\begin{moltablethree}
\hline
Data & Incidente & Consequência \\
\hline
Fev/1920 & Agressão a enfermeiro sem provocação & 30 dias de isolamento \\
Jun/1922 & \textbf{Homicídio --- Enfermeiro Wagner} & Isolamento por 90 dias \\
Mar/1937 & \textbf{Homicídio --- Paciente da ala C} & Isolamento permanente \\
Abr/1937 & Automutilação (arrancou unhas com os dentes) & Contenção física \\
Set/1945 & Episódio de vocalização incontrolável (alega que 'a Rasga Mortalha' está falando através dele) & Sedação com brometo \\
Jan/1948 & \textbf{Homicídio --- Visitante (entrou na ala proibida)} & Isolamento perpétuo \\
1949-1956 & Mutismo seletivo (7 anos sem falar com a equipe) & --- \\
1956 & Retomada da fala: primeira frase 'O diário' & --- \\
1979 & Sessões com Dr. Heitor Oliveira & Vide relatórios anexos \\
\hline
\end{moltablethree}

\vspace{0.5em}\noindent\rule{\textwidth}{0.5pt}\vspace{0.5em}

\noindent \textbf{Medicações Administradas (registro parcial):}

\noindent - Brometo de potássio (1920-1945) --- sedativo
- Eletrochoque (1937, 1945, 1952) --- sem anestesia
- Lobotomia transorbital --- \textit{indicada em 1948 mas não realizada (paciente reagiu com violência extrema à mera tentativa de aproximação)}
- Clorpromazina (1955 em diante) --- antipsicótico

\vspace{0.5em}\noindent\rule{\textwidth}{0.5pt}\vspace{0.5em}

\noindent \textbf{Observações Finais:}

\noindent Paciente 1.260 é um dos internos mais antigos do Colônia. Seu comportamento é imprevisível: períodos de calma absoluta intercalados com explosões de violência extrema. Mantém um diário que jamais permitiu que qualquer funcionário lesse em vida.

\noindent \textit{Nota adicional (1979): Dr. Heitor Oliveira obteve permissão para ler o diário. Primeira e única pessoa a ter acesso ao documento. Dr. Oliveira solicitou licença médica 48 horas depois. Não retornou ao trabalho.}

\noindent \textit{O diário encontra-se atualmente sob custódia do IML de Barbacena, anexado ao laudo de necropsia.}

\vspace{0.5em}
\noindent\textit{\textbf{↪ Rotas:} \rota{MEM-08} \rotasep \rota{DOC-03} \rotasep \rota{LUC-01}}
